\chapter{Reinforcement Learning Specification}
\label{app:C}

This appendix gives the complete specification of the RL-3 sequencing agent as implemented. It
supplements \Cref{sec:rl3}, which describes the design rationale.

\section{Agent architecture}

A single deep Q-network \citep{mnih2015} is shared across all three stages. The network maps a
16-dimensional state to two action values through one hidden layer of 64 units. The same
parameters are used at picking, packing and dispatch; the stage is supplied as an input feature
rather than by instantiating separate agents.

\begin{table}[htbp]
\centering
\caption{Network architecture.}
\label{tab:rl-network}
\small
\begin{tabularx}{0.85\textwidth}{lX}
\toprule
\textbf{Component} & \textbf{Specification} \\
\midrule
Input dimension  & 16 \\
Hidden layer     & 64 units \\
Output dimension & 2 (one action value per class) \\
Shared across    & picking, packing and dispatch \\
\bottomrule
\end{tabularx}
\end{table}

\section{State representation}

\Cref{tab:rl-state} lists the 16 features in the exact order in which they enter the network.
All features are clipped to \([0,1]\).

\begin{table}[htbp]
\centering
\caption[RL-3 state features]{The 16 state features, in implementation order.}
\label{tab:rl-state}
\scriptsize
\begin{tabularx}{\textwidth}{clp{2.6cm}X}
\toprule
\textbf{\#} & \textbf{Feature} & \textbf{Normalisation} & \textbf{Operational rationale} \\
\midrule
0  & Picking urgent queue   & \(/\,200\) & How much urgent work waits at picking \\
1  & Picking normal queue   & \(/\,500\) & How much normal work waits at picking \\
2  & Packing urgent queue   & \(/\,200\) & Urgent work waiting at packing \\
3  & Packing normal queue   & \(/\,500\) & Normal work waiting at packing \\
4  & Dispatch urgent queue  & \(/\,200\) & Urgent work waiting at dispatch \\
5  & Dispatch normal queue  & \(/\,500\) & Normal work waiting at dispatch \\
6  & Picking work in progress  & \(/\,5\) & Orders currently in service at picking \\
7  & Packing work in progress  & \(/\,5\) & Orders currently in service at packing \\
8  & Dispatch work in progress & \(/\,5\) & Orders currently in service at dispatch \\
9  & Elapsed time & \(/\,H\) & Position within the monthly horizon \\
10 & Urgent head slack & see below & Deadline pressure on the next urgent order here \\
11 & Normal head slack & see below & Deadline pressure on the next normal order here \\
12 & Stage identifier & \(0.0\,/\,0.5\,/\,1.0\) & Which stage the decision is being taken at \\
13 & Picking workers  & \(/\,20\) & Capacity available at picking \\
14 & Packing workers  & \(/\,20\) & Capacity available at packing \\
15 & Dispatch workers & \(/\,20\) & Capacity available at dispatch \\
\bottomrule
\end{tabularx}
\end{table}

Queue-length normalisers differ by class (\num{200} for urgent, \num{500} for normal) because
the normal class is by far the larger population; using a common divisor would compress the
urgent signal into a narrow range.

The slack features (10 and 11) describe the order at the head of each queue \emph{at the current
stage}:

\begin{equation}
\text{slack} \;=\; \frac{1}{2}\left(
\operatorname{clip}\!\left(\frac{\sigma_{c} - (\text{now} - a_i)}{\sigma_c},\, -1,\, 1\right) + 1
\right),
\label{eq:slack}
\end{equation}

\noindent where \(\sigma_c\) is the class's system-time SLA threshold. The expression maps
remaining time relative to the deadline onto \([0,1]\): values above
\num{0.5} indicate an order still within its deadline, values below indicate an order already
late. An empty queue yields the neutral value \num{0.5}.

The capacity features (13--15) are normalised by a fixed reference of 20 workers per stage
rather than by a per-episode maximum, so that a given worker count always maps to the same
feature value across configurations.

\paragraph{No lookahead.} None of the 16 features encodes future arrivals. The agent observes
present queue state, present work in progress, elapsed time, the urgency of the orders at the
head of each queue and the available capacity. It has no forecast and no privileged information
about what will arrive next.

\section{Action space and decision points}

The action space has two elements: action \num{0} serves the urgent queue, action \num{1} serves
the normal queue. Within the selected class, service is first-in-first-out.

The agent is consulted only when both queues at the stage are non-empty. If exactly one class is
waiting, the action is forced to that class and the transition is \emph{not} recorded for
training. If the network selects a class whose queue is empty, the action is overridden to the
non-empty class. Training therefore only sees decisions in which a genuine choice existed.

\section{Reward function}

The reward is deferred until the order completes dispatch, and is then applied to every
transition recorded for that order:

\begin{equation}
r_i =
\begin{cases}
w_{c(i)} & \text{on time} \\[4pt]
-\,p_{c(i)} \cdot \min\!\left(1, \dfrac{L_i}{\sigma_{c(i)} \cdot \kappa}\right) & \text{late}
\end{cases}
\qquad
r_i = -\,p_{c(i)} \ \ \text{if left as backlog at } H.
\label{eq:reward-app}
\end{equation}

\begin{table}[htbp]
\centering
\caption{Reward coefficients.}
\label{tab:rl-reward}
\small
\begin{tabularx}{0.9\textwidth}{llrX}
\toprule
\textbf{Symbol} & \textbf{Parameter} & \textbf{Value} & \textbf{Meaning} \\
\midrule
\(w_u\)      & On-time reward, urgent  & 5.0 & Value of meeting an urgent deadline \\
\(w_n\)      & On-time reward, normal  & 3.0 & Value of meeting a normal deadline \\
\(p_u\)      & Lateness penalty, urgent & 2.0 & Maximum penalty for a late urgent order \\
\(p_n\)      & Lateness penalty, normal & 2.0 & Maximum penalty for a late normal order \\
\(\kappa\)   & Lateness cap multiplier & 2.0 & Lateness saturates at twice the threshold \\
\bottomrule
\end{tabularx}
\end{table}

The on-time rewards differ by class (5.0 against 3.0), so meeting an urgent deadline is worth
more, while the lateness penalties are equal (2.0 each). The resulting swing between the on-time
and fully-late outcome is 7.0 for urgent orders against 5.0 for normal --- a ratio of
\num{1.4}. The agent therefore has an incentive to favour the urgent class without a structural
incentive to abandon the normal class.

Assigning the same terminal reward to every decision taken for an order is a credit-assignment
approximation. It is justified on the grounds that the service outcome is fully determined at
dispatch and every decision along the path contributed to it, but it does not distinguish
decisions that mattered from those that did not.

\section{Training configuration}

\begin{table}[htbp]
\centering
\caption{Training hyperparameters.}
\label{tab:rl-hyper}
\small
\begin{tabularx}{0.92\textwidth}{lrX}
\toprule
\textbf{Parameter} & \textbf{Value} & \textbf{Note} \\
\midrule
Episodes                & 200 & each samples one (month, configuration) pair \\
Learning rate           & \num{0.001} & Adam optimiser \\
Discount factor \(\gamma\) & 0.99 & \\
Batch size              & 256 & \\
Replay buffer capacity  & \num{200000} & \\
Training start          & \num{3000} transitions & minimum buffer before updates \\
Update frequency        & every 4 steps & \\
Updates per episode     & 500 & \\
Target network update   & every \num{2000} steps & \\
\(\varepsilon\) start / end & 1.0 / 0.05 & \\
\(\varepsilon\) decay steps & \num{200000} & \\
\bottomrule
\end{tabularx}
\end{table}

Each episode samples one (month, workforce configuration) pair uniformly from a fixed training
pool and runs the full month over its own operating horizon. Training deliberately omits the
common random number map described in \Cref{sec:crn}, so each episode encounters independent
service-time realisations. The episode's random seed is drawn from a generator initialised with
the simulation's base seed, so the sequence of sampled pairs and realisations is reproducible
given the same configuration.

\section{Training pool and holdout}
\label{sec:app-training-pool}

The training pool is constructed at the start of each training run and persisted as a manifest.
\Cref{tab:rl-pool} reproduces the manifest surviving in the repository. Three representative
months span the demand range; for each, the analytical estimate is computed and nine candidate
configurations are generated around it by the procedure of \Cref{sec:candidates}, of which the
last two are held out and the remaining seven form the training pool. Training therefore draws
from 21 (month, configuration) pairs, and six pairs are never trained on.

\begin{table}[htbp]
\centering
\caption[RL-3 training pool]{The persisted training-pool manifest
(\code{data/rl3\_train\_pool.json}).}
\label{tab:rl-pool}
\scriptsize
\begin{tabularx}{\textwidth}{llrXl}
\toprule
\textbf{Month} & \textbf{Role} & \textbf{Orders} & \textbf{Configurations} &
\textbf{Analytical centre} \\
\midrule
June & low & \num{9600} & \rgm{s432}, \rgm{s332}, \rgm{s532}, \rgm{s422}, \rgm{s442},
  \rgm{s431}, \rgm{s433} & (4,\,3,\,2) \\
 & \emph{holdout} & & \rgm{s321}, \rgm{s543} & \\
\midrule
October & medium & \num{20400} & \rgm{s964}, \rgm{s864}, \rgm{s10\_6\_4}, \rgm{s954},
  \rgm{s974}, \rgm{s963}, \rgm{s965} & (9,\,6,\,4) \\
 & \emph{holdout} & & \rgm{s753}, \rgm{s11\_7\_5} & \\
\midrule
December & peak & \num{40800} & \rgm{s26\_14\_7}, \rgm{s25\_14\_7}, \rgm{s27\_14\_7},
  \rgm{s26\_13\_7}, \rgm{s26\_15\_7}, \rgm{s26\_14\_6}, \rgm{s26\_14\_8} & (26,\,14,\,7) \\
 & \emph{holdout} & & \rgm{s22\_11\_5}, \rgm{s30\_17\_9} & \\
\bottomrule
\end{tabularx}
\end{table}

Three consequences of this design should be stated, because they bear directly on how the
results in \Cref{ch:results} generalise.

First, \textbf{nine of the twelve retrospective months were never trained on at all}. Only June,
October and December appear in the pool. Every result reported for the other nine months is an
out-of-distribution evaluation with respect to demand.

Second, \textbf{several of the configurations that carry the headline results are holdouts}. The
December configuration \rgm{s22\_11\_5} and the October configuration \rgm{s753} --- the two
cases examined most closely in \Cref{sec:res-policy} --- are both in the holdout set, as
confirmed independently by the persisted diagnostic report, which records
\code{seen\_during\_training: false} for them. That the agent performs well at those
configurations is therefore evidence of generalisation rather than of memorisation, which
strengthens the finding rather than qualifying it.

Third, \textbf{a second, unrelated train/holdout list exists in the configuration}, naming twelve
small static configurations for training and four (\rgm{s112}, \rgm{s231}, \rgm{s322},
\rgm{s432}) as holdouts. That list belongs to a separate generalisation-evaluation script and
predates the per-month dynamic candidate generation used here; it is \emph{not} what the trained
agent was trained on, and it is recorded here only to prevent the two from being confused.

\section{Evaluation mode}
\label{sec:app-greedy}

Every evaluation reported in this thesis invokes the network greedily: the action with the
highest estimated value is taken, with no random component, and no transition is recorded. The
\(\varepsilon\)-greedy schedule in \Cref{tab:rl-hyper} governs training only. This applies to the
retrospective grid, the forecast replications, the generalisation checks and the behavioural
diagnostics alike, so no reported RL-3 metric is affected by exploratory randomness. A
consequence worth noting is that RL-3's behaviour at a given configuration is deterministic given
the order stream and the service-time map; its results are not averaged over stochastic policy
draws.

\section{A representation limitation: feature saturation}
\label{sec:app-saturation}

All 16 features are clipped to \([0,1]\), which bounds the input scale but also means that any
underlying quantity exceeding its normalising constant maps to exactly \num{1.0}. Two groups of
features reach that ceiling in the experiments reported here.

The capacity features (13--15) divide the stage's worker count by \num{20}. Any stage staffed at
20 FTE or more therefore yields \num{1.0}, and the network cannot distinguish those capacities
through that feature. This is not hypothetical at peak: \emph{every one} of December's sixteen
candidate configurations staffs picking at 22 FTE or more --- from \rgm{s22\_11\_5} to
\rgm{s30\_17\_9} --- so the picking-capacity feature is identically \num{1.0} across the whole
December candidate set. Packing (11 to 17 FTE) and dispatch (5 to 9 FTE) remain below the
ceiling and continue to vary, as do the queue, work-in-progress, slack and elapsed-time features,
so the states are not identical; but one of the three capacity signals carries no information in
that month.

The queue-length features saturate in the same way. The normal-class divisor is \num{500}, while
the observed maximum picking queue in December exceeds \num{3000} orders and the mean picking
queue at \rgm{s22\_11\_5} is above \num{1000}. The persisted diagnostic report flags this
directly, recording a state-saturation signal at that configuration. During congested periods the
agent therefore sees a saturated queue signal that no longer conveys \emph{how} congested the
stage is.

Two things follow, and neither is that the policy is invalid. The clipped representation reduces
the network's resolution between sufficiently large capacity and queue values, so the agent
cannot condition finely on differences in that range; and because the saturating range coincides
with peak demand, the effect is concentrated in precisely the months where the sequencing
decision matters most. What the results show is that the agent nevertheless produced feasible,
lower-cost outcomes at those configurations, which suggests the remaining unsaturated features
carry enough signal. A representation with larger or adaptive normalising constants is a
straightforward extension, and is noted in \Cref{sec:future-work}.

\section{Checkpoint status}

The trained network weights were stored at \code{data/dqn\_rl3\_final.pt}. This path matches an
ignore rule in the repository configuration and was therefore never placed under version
control. At the time this thesis was prepared the file was \textbf{not present} in the working
tree.

Provenance metadata recorded during an earlier diagnostic session documents the file as
\num{24669} bytes with SHA-256 digest beginning \code{2002aaab...0f2eed}, and records that it
loaded into the current 16--64--2 architecture without shape mismatch across repeated evaluation
runs. That metadata is retained as evidence of what was evaluated: it identifies the artefact
and establishes that it matched the architecture documented here. It does \emph{not} permit the
artefact to be reconstructed --- a digest verifies a file one already holds, and cannot produce
one.

The consequence for this thesis is stated plainly. The persisted evaluation outputs are
available and every RL-3 figure and table derives from them, but exact regeneration of the RL-3
results requires the trained checkpoint, which was not available in the environment used for
final thesis preparation. Retraining from the recorded configuration would produce \emph{a}
policy, not \emph{this} policy, since training samples episodes stochastically and the resulting
weights would differ. No part of this thesis claims complete end-to-end reproducibility of the
RL-3 results from the distributed artefacts. This is a genuine limitation affecting one of the
three compared policies, and is recorded as such in \Cref{sec:disc-limitations} rather than
minimised.
